\subsection{Cache và quản lý độ mới}
\label{subsec:db_cache}

Cache được kiểm soát bằng ba cơ chế. Cơ chế thứ nhất là TTL dựa trên thời điểm cập nhật. Dữ liệu overview, market hoặc kết quả phân tích được trả trực tiếp khi vẫn còn mới; khi hết TTL, service gọi provider và cập nhật lại bản ghi.

Cơ chế thứ hai là coverage metadata. Với lịch sử ví và enhanced transaction, metadata lưu khoảng thời gian đã đồng bộ. Service so sánh khoảng yêu cầu với khoảng đã có, chỉ tải phần còn thiếu rồi hợp nhất coverage. Cách này tránh tải lại toàn bộ lịch sử của ví.

Cơ chế thứ ba là upsert. Market snapshot, chart point và các cache theo khóa nghiệp vụ sử dụng insert kết hợp conflict handling. Bản ghi mới được thêm khi khóa chưa tồn tại; dữ liệu cũ được cập nhật khi cùng address, timestamp hoặc signature.

\subsection{Toàn vẹn và bảo mật dữ liệu}
\label{subsec:db_toan_ven}

Primary key tổng hợp được dùng cho dữ liệu chuỗi thời gian và lịch sử, ví dụ address kết hợp timestamp hoặc transaction hash. Unique constraint ngăn một transaction detail hoặc chart point bị ghi nhiều lần. Foreign key với quy tắc cascade áp dụng cho dữ liệu thuộc người dùng như auth account, watchlist, alert và subscription.

Không phải mọi liên hệ logic đều được chuyển thành foreign key. Wallet address và token address đến từ blockchain, có thể xuất hiện trước khi metadata tương ứng được tải. Việc bắt buộc foreign key trong trường hợp này sẽ làm gián đoạn luồng đồng bộ. Service chịu trách nhiệm chuẩn hóa address và giữ quy tắc ghép dữ liệu nhất quán.

Dữ liệu công khai từ blockchain được tách khỏi dữ liệu cá nhân. Các truy vấn profile, watchlist, cảnh báo, chat và thanh toán phải được giới hạn theo user ID đã xác thực. Password chỉ được lưu dưới dạng hash; API key và secret của provider không được lưu trong các bảng nghiệp vụ.

% TODO(runtime): Hoàn thiện ghi alert_history sau khi delivery thành công.
% TODO(runtime): Bổ sung API lấy lịch sử, đánh dấu đã đọc và chống gửi trùng.

\subsection{Tổng kết thiết kế cơ sở dữ liệu}
\label{subsec:db_tong_ket}

Thiết kế cơ sở dữ liệu kết hợp dữ liệu quan hệ với các read model cache theo miền nghiệp vụ. Các bảng user, payment và alert ưu tiên ràng buộc quan hệ. Các bảng token và wallet ưu tiên khả năng đồng bộ, truy vấn theo address và tái sử dụng dữ liệu provider. Việc tách overview, history, analysis và transaction detail giúp mỗi endpoint đọc đúng cấu trúc cần thiết thay vì phụ thuộc vào một bảng giao dịch chung.
